\documentclass[12pt,a4paper]{article}
\usepackage[utf8]{inputenc}
\usepackage{amsmath, amssymb, amsfonts}
\usepackage{geometry}
\geometry{margin=2.5cm}
\usepackage{hyperref}
\usepackage{pgfplots}
\usepackage{fancyhdr}
\usepackage{listing}
\usepackage{listings}
\usepackage{titlesec}
\titleformat{\section}
  {\normalfont\Large\bfseries\centering}  % Format: font size, bold, centered
  {\thesection}{1em}{}
\pagestyle{fancy}
\pgfplotsset{compat=1.18}
\fancyhead{}
\fancyfoot{}
\title{Laboratorio 2}
\author{}
\date{}
\fancyhead[L]{Gabriele Perna}          % Left side
\fancyhead[C]{Laboratorio 2}     % Center
\fancyhead[R]{2026/2027}             % Right side
\fancyfoot[C]{\thepage}           % Page number in the center of footer
\let\oldsection\section
\renewcommand{\section}{\clearpage\oldsection}

\begin{document}

\maketitle
\tableofcontents

\section{Introduzione}
In questo corso lavoreremo con programmi vicini al linguaggio macchina.\\
Lavoreremo inoltre con:
\begin{itemize}
    \item UNIX
    \item thread
    \item processi
    \item programmazione di sistema 
\end{itemize}
\paragraph{Codice Teams:} r407rxv

\section{Il terminale UNIX}
Il \textbf{terminale} è una interfaccia testuale per dialogare con il sistema.\\
Dentro il terminale gira una \textbf{shell}, ovvero il programma che interpreta i comandi.\\
Il terminale mostra il \textbf{prompt}, riceve i comandi e visualizza l'output.\\
Il simbolo $\$$ indica che la shell è pronta a ricevere un comando (potrebbe non essere $\$$ a volte).
\subsection{Comandi}
\subsubsection{Struttura base comandi Linux}
Esiste una struttura base per molti comandi, ovvero questa:
\begin{center}
    comando + opzione + argomento + (opzione + argomento)
\end{center}
La differenza principale tra opzione e argomento è che l'opzione è relativa alla \textbf{modalità} del comando e l'argomento \textbf{su cosa lavora} il comando. L'opzione + argomento invece è relativo all'argomento.
per esempio:
\begin{lstlisting}
    $ gcc -Wall hello.c -o hello
\end{lstlisting}
\subsubsection{--help e man}
Entrambi questi comandi aiutano a comprendere specifici comandi in modi diversi, --help è un aiuto generico mentre "man" è il manuale.
\begin{lstlisting}
    $ [comando] --help
    $ gcc --help
    $ man [comando]
    $ man gcc
\end{lstlisting}
\subsubsection{pwd}
\textbf{pwd} è il comando "print working directory" e mostra la posizione della directory corrente all'interno del \textbf{file system}.
\begin{lstlisting}
    $ pwd
\end{lstlisting}
\subsubsection{ls}
\textbf{ls}, ovvero "list", elenca tutto il contenuto della directory corrente, directory e file.
\begin{lstlisting}
    $ ls
\end{lstlisting}
Può essere usata con alcuni \textbf{attributi}, aggiungibili con un "dash" (-) come in questo caso:
\begin{lstlisting}
    $ ls -a
\end{lstlisting}
Ecco gli attributi:
\begin{itemize}
    \item \textbf{a}: include tutti gli elementi nascosti
    \item \textbf{l}: include informazioni avanzate (data, grandezza...)
\end{itemize}
\subsubsection{cd}
\textbf{cd} significa "change directory" e si sposta nella directory specificata da un path relativo o assoluto.
\begin{lstlisting}
    $ cd [path]
    $ cd esercizi
    $ cd /users/gabri/esercizi
\end{lstlisting}
Percorsi speciali:
\begin{itemize}
    \item \textbf{Directory corrente}: specificato da un punto, ha lo stesso valore di pwd.
    \begin{lstlisting}
       $ cd . 
    \end{lstlisting}
    \item \textbf{Directory padre}: specificato da un doppio punto.
    \begin{lstlisting}
       $ cd .. 
    \end{lstlisting}
    \item \textbf{Home dell'utente}: specificato da una tilde, si sposta nella home dell'utente ovunque essa sia.
    \begin{lstlisting}
        $ cd ~
    \end{lstlisting}
\end{itemize}
\subsubsection{mkdir}
\textbf{mkdir} significa "make directory" e crea una directory del nome specificato nella cartella corrente.
\begin{lstlisting}
    $ mkdir [nomeDirectory]
    $ mkdir esercizi
\end{lstlisting}
\subsubsection{touch}
Il comando \textbf{touch} nasce con l'intento di aggiornare l'ultimo accesso e l'ultima modifica di un file (marcatura temporale, timestamp).
\begin{lstlisting}
    $ touch [nomeFile]
    $ touch log.txt
\end{lstlisting}
Se il file non esiste, lo crea.\\
Tale comando senza attributi modifica:
\begin{itemize}
    \item Tempo di accesso (atime)
    \item Tempo di modifica (mtime)
\end{itemize}
Se desideriamo poi modificare uno solo di questi timestamp allora si può introdurre un attributo:
\begin{itemize}
    \item \textbf{-a}: atime
    \item \textbf{-m}: mtime
\end{itemize}
Il comando può essere usato anche sulle directory, ma non crea automaticamente una directory se assente, bensì solo file.
\subsubsection{stat}
Il comando \textbf{stat}, ovvero "status", mostra i metadati di un file o una directory, inclusi i timestamp.
\begin{lstlisting}
    $ stat [nome]
    $ stat hello.c
\end{lstlisting}
Ecco i dati mostrati da tale comando:
\begin{itemize}
    \item Tipo di file
    \item Dimensione
    \item Inode
    \item Permessi
    \item Proprietario e Gruppo
    \item Data di ultimo accesso (Access)
    \item Data di ultima modifica (Modify)
    \item Data di ultima modifica dei metadati (Change)
\end{itemize}
\subsubsection{cp}
Il comando \textbf{cp}, ovvero "copy", copia il file passato come parametro e lo incolla nel path specificato dopo l'argomento, eventualmente rinominandolo.
\begin{lstlisting}
    $ cp [nomeFileDaCopiare] [Destinazione]\[NuovoNomeFile]
    $ cp hello.c \esercizi\eserc.c
\end{lstlisting}
Questo esempio copia il file "hello.c", lo sposta nella sub-directory "esercizi" e lo rinomina.\\
Il file originale rimane al suo posto.\\
La rinomina della copia è omissibile.\\
Questo comando, con l'aggiunta di un attributo "-r" (\textbf{recursive}) può essere applicato anche alle directory.
\begin{lstlisting}
    $ cp -r [nomeDirectory] [Destinazione]\[NuovoNomeDirectory]
    $ cp -r cartella1 \esercizi\cartella2
\end{lstlisting}
\subsubsection{mv}
Il comando \textbf{mv}, ovvero "move", sposta o/e rinomina un file:
\begin{lstlisting}
    $ mv [nomeVecchio] [destinazione]\[nomeNuovo]
    $ mv hello.c ciao.c
\end{lstlisting}
Questo comando rinomina il file "hello.c" in "ciao.c", senza spostarlo.
\subsubsection{rm}
Il comando \textbf{rm}, ovvero "remove", elimina un file.
\begin{lstlisting}
    $ rm [nomeFile]
    $ rm hello.c
\end{lstlisting}
La cancellazione da terminale non sposta nel cestino, è in genere definitiva.
\subsubsection{rmdir}
Il comando \textbf{rmdir}, ovvero "remove directory", elimina una directory vuota.
\begin{lstlisting}
    $ rmdir [nomeDirectory]
    $ rmdir esercizi
\end{lstlisting}
Se vogliamo eliminare una cartella e il suo contenuto, si possono usare questi attributi:
\begin{itemize}
    \item \textbf{-r}: "recursive", applica la cancellazione anche ai contenuti.
    \item \textbf{-f}: "forced", applica l'attributo recursive con forza, ignora le conferme di cancellazione.
\end{itemize}
\subsubsection{cat}
Il comando \textbf{cat}, ovvero "concatenate", scrive il contenuto del file sullo standard output.
\begin{lstlisting}
    $ cat [nomeFile]
    $ cat hello.txt
\end{lstlisting}
Si possono stampare più file in sequenza, concatenati:
\begin{lstlisting}
    $ cat [nomeFile1] [nomeFile2]
    $ cat hello1.txt hello2.txt
\end{lstlisting}
Il comando usato da solo diventa interattivo e legge dallo standard input, ovvero il terminale: ogni volta che si scrive qualcosa a cat, lo stamperà direttamente sul terminale.
\begin{lstlisting}
    $ cat
    $ ciao
    ciao
\end{lstlisting}
Per uscire da questa modalità bisogna premere (\textbf{Ctrl + D}).
\subsubsection{less}
Il comando \textbf{less}, opera simile al comando \textbf{cat}, ma non carica tutto il file nel terminale e per leggere più contenuto bisogna scrollare.
\begin{lstlisting}
    $ less [nomeFile]
    $ less hello.txt
\end{lstlisting}
Si possono aprire anche più file insieme:
\begin{lstlisting}
    $ less [nomeFile1] [nomeFile2]
    $ less file1.txt file2.txt
\end{lstlisting}
Per passare al prossimo file o al precedente si fa rispettivamente:
\begin{lstlisting}
    $ :n
    $ :p
\end{lstlisting}
Per scorrere in avanti si può premere:
\begin{itemize}
    \item \textbf{Spazio}: pagina successiva.
    \item \textbf{$\downarrow$}: riga successiva.
    \item \textbf{Invio/Enter}: riga successiva.
    \item \textbf{G}: fine del file.
\end{itemize}
Questo comando permette anche di scorrere indietro:
\begin{itemize}
    \item \textbf{b}: pagina precedente
    \item \textbf{Freccia in su}: riga precedente
    \item \textbf{g}: inizio del file
\end{itemize}
e di cercare parole, vediamo un esempio (una volta attivato less):
\begin{lstlisting}
    $ /ciao (cerca in avanti)
    $ ?ciao (cerca indietro)
\end{lstlisting}
Per scorrere le occorrenze si fa cosi:
\begin{lstlisting}
    $ n (occorrenza successiva)
    $ N (occorrenza precedente)
\end{lstlisting}
Per uscire bisogna premere \textbf{Q}.
\subsubsection{more}
Il comando \textbf{more} funziona come il comando \textbf{less}, ma non permette di scorrere indietro, nè di cercare parole.
\begin{lstlisting}
    $ more [nomeFile]
    $ more file.txt
\end{lstlisting}
Per uscire bisogna premere \textbf{Q}.\\
\subsubsection{head}
Il comando \textbf{head} permette di stampare le prime righe di un file su terminale.
\begin{lstlisting}
    $ head [nomeFile]
    $ head hello.txt
\end{lstlisting}
Possiamo specificare quante righe stampare:
\begin{lstlisting}
    $ head -n [righeDaStampare] [nomeFile]
    $ head -n 5 hello.txt
\end{lstlisting}
Questo esempio stampa le prime 5 righe del file "hello.txt".
\subsubsection{tail}
Il comando \textbf{tail} permette di stampare le ultime righe di un file su terminale.
\begin{lstlisting}
    $ tail [nomeFile]
    $ tail hello.txt
\end{lstlisting}
Possiamo specificare quante righe stampare:
\begin{lstlisting}
    $ tail -n [righeDaStampare] [nomeFile]
    $ tail -n 5 hello.txt
\end{lstlisting}
Questo esempio stampa le ultime 5 righe del file "hello.txt".
\subsubsection{wc}
Il comando \textbf{wc}, ovvero "word count", serve per contare delle specifiche informazioni in un file.
\begin{lstlisting}
    $ wc [nomeFile]
    $ wc file.txt
\end{lstlisting}
Il comando restituisce 3 valori in questo ordine:
\begin{itemize}
    \item Numero di righe
    \item Numero di parole
    \item Numero di byte
\end{itemize}
Il comando può essere eseguito con un filtro passato come attributo:
\begin{itemize}
    \item \textbf{-l} per le righe
    \item \textbf{-w} per le parole
    \item \textbf{-c} per i byte
\end{itemize}

\subsection{Wildcards}
Le wildcard permettono di mantenere un aspetto generico di alcuni dati per lavorare su più entità in contemporanea.\\
Le wildcard quindi funzionano da veri e propri \textbf{placeholder}.
Questa è una funzionalità della shell, la quale \textbf{espande} le wildcard una volta che gli arrivano: sostituisce le wildcard e genera un elenco di tutte le unità che soddisfano la wildcard.\\
Per esempio:
\begin{lstlisting}
    $ ls *.c
    "espansione shell"
    $ ls file1.c file2.c file3.c
\end{lstlisting}
Ai comandi non viene mai passata la shell, bensì gli argomenti precisi su cui lavorare.
\subsubsection{Asterisco (*)}
L'asterisco (*) usato come wildcard permette di specificare una \textbf{stringa} \textbf{qualsiasi} di lunghezza qualsiasi, per esempio:
\begin{lstlisting}
    $ ls *.c
\end{lstlisting}
Questo comando stampa tutti i file con estensione ".c", ovvero i file con nome qualsiasi che hanno estensione ".c".
\subsubsection{Punto Interrogativo (?)}
Il punto interrogativo (?) usato come wildcard indica un \textbf{carattere} qualsiasi.
\begin{lstlisting}
    $ ls prova?.c
\end{lstlisting}
Questo comando stampa tutti i file con prefisso "prova", suffisso ".c" e un carattere qualsiasi che divide suffisso e prefisso specificati, in questo caso utile per stampare tutte le prove / test possibilmente nominati "prova1", "prova2" e cosi via.
\subsection{Reindirizzazione}
\subsubsection{Reindirizzare l'output}
Si può reindirizzare lo \textbf{standard output} di un comando in questo modo:
\begin{lstlisting}
    $ [comando] > [destinazione]
    $ ls > lsLog.txt
\end{lstlisting}
Questo comando scrive in "lsLog.txt" quello che sarebe dovuto essere stampato sul terminale dal comando "ls".
\paragraph{Attenzione:} questo tipo di reindirizzazione \textbf{sovrascrive} ogni volta il file intero, se si desidera eseguire un'aggiunta in fondo al file senza perdere i dati precedenti allora bisogna usare l'\textbf{append}:
\begin{lstlisting}
    $ [comando] >> [destinazione]
    $ ls >> lsLog.txt
\end{lstlisting}
\subsubsection{Reindirizzare l'input}
Come l'output, si può reindirizzare lo \textbf{standard input}:
\begin{lstlisting}
    $ [destinazione] < [sorgente]
    $ sort < dati.txt
\end{lstlisting}
In questo esempio, si stampa in ordine alfabetico il contenuto di "dati.txt"
\subsubsection{Reindirizzare gli errori}
Anche il flusso \textbf{standard error} può essere reindirizzato:
\begin{lstlisting}
    $ [comando] 2> [destinazione]
    $ ls 2> lsLog.txt
\end{lstlisting}
\subsection{Pipe}
La \textbf{pipe} collega più comandi in questo modo:
\begin{center}
    stdout del primo comando $\rightarrow$ stdin del secondo comando
\end{center}
Ecco il template e un esempio:
\begin{lstlisting}
    $ [comando1] | [comando2]
    $ ls | less
\end{lstlisting}
In questo esempio viene eseguito "ls" e viene poi stampato secondo il comando "less".\\
La pipe evita di creare file intermedi che sarebbero necessari per realizzare il solito effetto.
\subsection{Virgolette}
Per evitare che non ci sia ambiguità negli argomenti, dato che la shell utilizza gli spazi per separare argomenti, comandi e opzioni, bisogna utilizzare le \textbf{virgolette}:
\begin{lstlisting}
    $ mkdir "nuova cartellozza"
\end{lstlisting}

\section{Filesystem UNIX}
Il \textbf{filesystem} è un albero che ha come radice la \textbf{root}, indicata da un "/"
\subsection{Path}
Il filesystem localizza ogni file e directory tramite un \textbf{path}, il quale può essere:
\begin{itemize}
    \item \textbf{Assoluto:} parte dalla root.
    \item \textbf{Relativo:} parte dalla directory corrente.
\end{itemize}
\subsection{Marcature temporali}
tenere traccia del tempo in UNIX è importante: i timestamp sono spesso utilizzati per prendere decisioni importanti, come la ricompilazione di un file con un timestamp vecchio.\\
In UNIX ogni file ha 3 marcature temporali:
\begin{itemize}
    \item \textbf{atime}: access time, ultimo accesso.
    \item \textbf{mtime}: modification time, ultima modifica del contenuto.
    \item \textbf{ctime}: change time, ultima modifica dei metadati.
\end{itemize}
\subsection{File nascosti}
I \textbf{file nascosti} sono i file il cui nome solitamente inizia con il punto.
\begin{center}
    .txt
    .config
    .bashrc
\end{center}

\section{Programmi UNIX}
\subsection{Flussi standard}
Ogni programma ha 3 flussi standard:
\begin{itemize}
    \item \textbf{standard input}
    \item \textbf{standard output}
    \item \textbf{standard error}
\end{itemize}
Questi flussi possono essere reindirizzati altrove dal terminale.

\section{Linguaggio C}
Il linguaggio \textbf{C} è un linguaggio compilato, pertanto, prima di eseguire un file bisogna convertirlo in binario.\\
C è un linguaggio \textbf{strongly typed}, c'è un forte controllo dei tipi.\\
Necessita ";" (semicolon) alla fine di ogni istruzione.\\
Esistono diversi standard C: C89, C95, C99, C11, C17, C23.
\subsection{Struttura generica}
Un programma C solitamente è composto da questi elementi in questo ordine:
\begin{enumerate}
    \item Direttive pre-processore
    \item Dichiarazioni globali
    \item Funzioni (incluso main)
\end{enumerate}
\paragraph{ATTENZIONE:} un programma C può avere \textbf{un solo main}.
\subsection{Compilatore}
Il compilatore più conosciuto e usato per C è \textbf{gcc}, ovvero "GNU C Compiler", anche se adesso viene definito come "GNU Compiler Collection" perche diventato un progetto più ampio che include anche altri compilatori come quello di Fortran e C++. 
\begin{lstlisting}
    gcc hello.c -o hello
\end{lstlisting}
In questo caso, il comando:
\begin{itemize}
    \item compila il file "hello.c"
    \item rinomina l'eseguibile come "hello"
\end{itemize}
Ogni programma eseguito ritorna un risultato per determinare se è terminato con successo o meno, si ottiene cosi:
\begin{lstlisting}
    $ gcc hello.c
    $ hello.c
    $ echo $?
    0
\end{lstlisting}
\subsubsection{Warnings}
I \textbf{warnings} sono importanti e non vanno ignorati, per compilare e stamparli tutti si fa cosi:
\begin{lstlisting}
    $ gcc -Wall hello.c
\end{lstlisting}
Ogni warning contiene un messaggio e la location (riga e carattere) di esso.
\subsection{Fasi di pre-esecuzione}
Ecco cosa succede dall'avvio del programma:
\begin{itemize}
    \item \textbf{Preprocessing}: vengono eseguite le direttive pre-processore.
    \item \textbf{Compilazione}: il sorgente .c viene trasformato in codice oggetto .o.
    \item \textbf{Linking}: i file oggetto vengono collegati in un unico eseguibile.
    \item \textbf{Caricamento}: l'eseguibile viene caricato in memoria per essere eseguito, viene diviso in 4 \textbf{zone} diverse:
    \begin{itemize}
        \item \textbf{Codice}: contiene le istruzioni e le funzioni.
        \item \textbf{Dati statici}: variabili globali, dati che rimangono per tutta la durata dell'esecuzione.
        \item \textbf{Heap}: zona dinamica, mantiene i dati temporanei come variabili ed oggetti.
        \item \textbf{Stack}: zona dinamica, quando una funzione entra in esecuzione viene generato un \textbf{record di attivazione}, il quale contiene varie informazioni legate a tale chiamata. Uno stack viene eliminato al terminare di un blocco di funzione.
    \end{itemize}
\end{itemize}
Vediamo come svolgere queste fasi manualmente invece che compilare tutto insieme direttamente:
\begin{itemize}
    \item Preprocessing: con "-E" ci si ferma alle direttive di preprocessing, "-o" invece salva l'output in un file secondario.
    \begin{lstlisting}
        $ gcc -E [nomeSorgente] -o [nuovoNome]
        $ gcc -E codice.c -o codiceP.c
    \end{lstlisting}
    \item Compilazione: si accettano solo file ".o", qui con "-c" si compila ma ci si ferma prima del Linking.
    \begin{lstlisting}
        $ gcc -c [nomeFile] -o [nuovoNome]
        $ gcc -c codiceP.c -o codice.o
    \end{lstlisting}
    \item Linking:
    \begin{lstlisting}
        $ gcc [nomeFile] -o [nuovoNome]
        $ gcc codice.o -o eseguibile.out
    \end{lstlisting}
\end{itemize}

\subsection{Librerie}
Nei programmi in C è possibile utilizzare delle \textbf{direttive di preprocessazione} (o del preprocessore) che permettono di importare funzioni e strumenti utili allo svolgimento del programma:
\begin{lstlisting}
    #include <stdio.h>
\end{lstlisting}
\subsection{Variabili}
In C possiamo usare le \textbf{variabili}. Sono celle di memoria a cui si può associare un valore di un tipo specifico, assegnato al momento della \textbf{dichiarazione}.
\begin{lstlisting}
    int x; //Dichiarazione
    x = 32; //Assegnamento
\end{lstlisting}
La dichiarazioni servono al compilatore per capire quanta memoria riservare.\\
Le variabili possono essere di molti tipi:
\begin{itemize}
    \item \texttt{char}: carattere, 1 byte
    \item \texttt{short}: numero intero, almeno 2 byte (piu corto di int comunque)
    \item \texttt{int}: numero intero, almeno 2 byte
    \item \texttt{long}: numero intero, almeno 4 byte
    \item \texttt{long long}: numero intero, almeno 8 byte
    \item \texttt{float}: numero reale a precisione singola, almeno 4 byte
    \item \texttt{double}: numero reale a precisione doppia, almeno 8 byte
    \item \texttt{long double}: numero reale a precisione estesa, almeno 10 byte
    \item \texttt{void}: assenza di tipo, 0 byte
\end{itemize}
La presenza della parola "almeno" consiste nel fatto che non si sa mai quanto veramente sia grande una variabile e in quanti bit consista, però alcune parole chiave garantiscono una lunghezza minima.\\
I char possono essere assegnati a numeri corrispondenti a caratteri ASCII.\\
Come si può vedere non esiste nessun tipo \textbf{string}, bensì sono rappresentate come array di char:
\begin{lstlisting}
    char y [8] = "Gabriele";
\end{lstlisting}
L'insieme dei valori delle variabili di un programma C in un determinato momento rappresenta lo \textbf{stato} del programma.
\subsubsection{Scoping}
Quando dichiariamo una variabile dobbiamo sapere quando inizia ad essere visibile e quando non lo è più, per questo si deve studiare lo \textbf{scoping}. Ci sono due tipi di scoping:
\begin{itemize}
    \item \textbf{Scoping globale}: tutte le variabili dichiarate fuori dalle funzioni (anche main)
    \item \textbf{Scoping locale}: le variabili definite dentro un blocco sono visibili solo nel blocco stesso.
\end{itemize}
\subsection{printf}
Questo è un comando base per stampare qualcosa su console:
\begin{lstlisting}
    printf("Hello World");
\end{lstlisting}
\subsubsection{Placeholders}
Per stampare le variabili bisogna utilizzare i \textbf{placeholders}:
\begin{lstlisting}
    int x = 10;
    printf("%d",x)
\end{lstlisting}
Bisogna scegliere il placeholder corretto per riuscire a stampare una variabile; la scelta va in base al tipo:
\begin{itemize}
    \item \textbf{\%d}: numeri interi (int)
    \item \textbf{\%i}: numeri interi (int)
    \item \textbf{\%f}: numeri decimali (double)
    \begin{itemize}
        \item Per stampare uno specifico numero di cifre decimali si aggiunge un "." dopo "\%" e poi si specifica il numero:
        \begin{lstlisting}
            printf(%.2f,x)
        \end{lstlisting}
        Questa stampa mostrerà per esempio solo 2 cifre decimali. Usato sugli altri placeholder questo provoca l'ingrandimento di padding a destra (numero negativo) o sinistra (numero positivo) di ciò che viene stampato.
    \end{itemize}
    \item \textbf{\%c}: caratteri (char)
    \item \textbf{\%s}: stringa
    \item \textbf{\%u}: numeri interi senza segno (unsigned int)
    \item \textbf{\%x}: numero intero in esadecimale
    \item \textbf{\%p}: indirizzo di memoria, puntatore
    \item \textbf{\%z}: tipo di sizeof (size\_t)
\end{itemize}
\subsection{Costanti}
In C non esistono le variabili costanti, possono essere sempre modificate!
\subsection{Commenti}
I commenti in C possono essere inseriti in vari modi:
\begin{itemize}
    \item Commenti su una riga: aperti da "//", non vengono chiusi
    \item Commenti su più righe: delimitati da "/*" e da "*/".
\end{itemize}
I commenti sono fondamentali e spariscono al momento di compilazione.\\
Sarebbe preferibile utilizzare sempre i commenti con "/*" perchè quelli con "//" sono più idonei per C++.
\subsection{Operatori in C}
\textbf{Operatori aritmetici}
\begin{itemize}
    \item \texttt{+}: addizione
    \item \texttt{-}: sottrazione
    \item \texttt{*}: moltiplicazione
    \item \texttt{/}: divisione
    \item \texttt{\%}: resto della divisione (modulo)
\end{itemize}
Bisogna stare attenti perchè alcune operazioni possono ritornare errori pratici di \textbf{overflow} o \textbf{underflow}.
\textbf{Operatori logici}
\begin{itemize}
    \item \texttt{\&\&}: AND logico
    \item \texttt{||}: OR logico
    \item \texttt{!}: NOT logico
\end{itemize}
\textbf{Operatori di assegnamento}
\begin{itemize}
    \item \texttt{=}: assegnamento
    \item \texttt{+=}: somma e assegnamento
    \item \texttt{-=}: sottrazione e assegnamento
    \item \texttt{*=}: moltiplicazione e assegnamento
    \item \texttt{/=}: divisione e assegnamento
    \item \texttt{\%=}: modulo e assegnamento
\end{itemize}
L'assegnamento avviene da destra verso sinistra:
\begin{lstlisting}
    a = b = c; ---> a = (b = c)
\end{lstlisting}
\textbf{Operatori di incremento e decremento}
\begin{itemize}
    \item \texttt{++}: incremento di 1
    \item \texttt{--}: decremento di 1
    \item \texttt{++x}: pre-incremento (incrementa e poi assegna)
    \item \texttt{x++}: post-incremento (assegna e poi incrementa)
    \item \texttt{--x}: pre-decremento (decrementa e poi assegna)
    \item \texttt{x--}: post-decremento (assegna e poi decrementa)
\end{itemize}
\textbf{Operatori di confronto}
\begin{itemize}
    \item \texttt{==}: uguale
    \item \texttt{!=}: diverso
    \item \texttt{>}: maggiore
    \item \texttt{<}: minore
    \item \texttt{>=}: maggiore o uguale
    \item \texttt{<=}: minore o uguale
\end{itemize}
\subsection{sizeof}
La keyword \textbf{sizeof} restituisce quanti byte occupa in memoria un tipo di dato o una variabile:
\begin{lstlisting}
    int x = 10;
    printf("%zu\n",sizeof(x))
\end{lstlisting}
Il placeholder utilizzato (z) è specificato in una sottosezione precedente.
\subsection{If-Else}
Per prendere decisioni si usano i costrutti \textbf{if-else}:
\begin{lstlisting}
    if (condizione) {
        comando;
        ...
    } else {
        comando;
        ...
    }
\end{lstlisting}
La condizione deve essere paragonabile con un valore \textbf{booleano} per determinare la decisione.\\
Questi costrutti si possono annidare e addirittura si può omettere l'else.
\subsection{While e do-while}
Questo costrutto permette di eseguire più volte le stesse istruzioni fino ad un certo \textbf{break point}, ovvero fino a che una condizione non ferma il while.
\begin{lstlisting}
    while (condizione) {
        comando;
        ...
    }
\end{lstlisting}
C'è un'opzione alternativa che garantisce una prima esecuzione e avvia il ciclo solo se la condizione è soddisfatta:
\begin{lstlisting}
    do {
        comando;
        ...
    } while (condizione)
\end{lstlisting}

\section{Esercizi}
\subsection{Esercizio 1}
Partendo dalla propria home, svolgere queste operazioni dal terminale LINUX:
\begin{enumerate}
    \item Creare una directory lezione02
    \item Entrarci
    \item Creare tre file vuoti: "a.c", "b.c", "note.txt"
    \item Visualizzare solamente i file con estensione ".c"
    \item Salvare quell'elenco in "sorgenti.txt"
    \item Aggiungere il contenuto di "sorgenti.txt" a "note.txt"
    \item Contare quante righe contiene "note.txt"
\end{enumerate}
\textbf{Soluzione:}
\begin{lstlisting}
    $ cd ~
    $ mkdir lezione02
    $ cd lezione02
    $ touch a.c b.c node.txt
    $ ls *.c > sorgenti.txt
    $ sorgenti.txt >> note.txt
    $ qw -l note.txt
\end{lstlisting}

\end{document}